\title{LongRoPE: Extending LLM Context Window Beyond 2 Million Tokens}

\begin{abstract}
A large context window is a valuable capability for large language models (LLMs). Nevertheless, present extended context windows generally max out at approximately 128k tokens because of the substantial expenses associated with fine-tuning, limited availability of lengthy textual data, and the disruptive effects caused by novel token positions.
\end{abstract}

This work presents LongRoPE, a novel approach that expands the context window of pre-trained LLMs to an unprecedented \textbf{2048k} tokens, using no more than 1k fine-tuning steps within a 256k token training length, all while preserving the model’s original performance on short contexts. This advancement is grounded in three primary contributions: (i) by uncovering and utilizing two distinct types of positional interpolation non-uniformities through an optimized search process, we provide improved fine-tuning initialization and facilitate an 8$\times$ extension even without additional fine-tuning; (ii) we propose a stepwise extension framework, beginning with fine-tuning at 256k tokens and then applying a second round of positional interpolation on the already fine-tuned, length-extended LLM to ultimately support a 2048k token context window; (iii) we recalibrate LongRoPE at an 8k window size to restore performance on short context tasks. Through thorough experimentation on LLaMA2 and Mistral over a broad spectrum of benchmarks, we validate the efficacy of LongRoPE. The approach preserves the original model architecture, requiring only slight adjustments to the positional embedding mechanism, and is compatible with most existing optimizations. The codebase will be released at \texttt{https://github.com/microsoft/LongRoPE}

\section{Introduction}

Although Large Language Models (LLMs) have achieved significant progress across a range of tasks~\cite{gpt4,llama2}, they are still constrained by a finite context window, such as the 4096-token limitation seen in LLaMA2~\cite{llama2}. When input exceeds this context length, the model’s effectiveness deteriorates, as it encounters positional configurations absent from its training. This limitation introduces difficulties for key applications, including in-context learning with a large number of exemplars~\cite{Huang2023FewerIM} as well as in the deployment of LLM agents~\cite{park2023generative,madaan2023selfrefine}.

Recent studies demonstrate that the context window of a pre-trained LLM can be enlarged to approximately 128k through fine-tuning with longer input sequences~\cite{longlora,pi,yarn,zhang2024soaring,liu2023scaling}. However, further increasing the context window faces three main challenges. First, introducing position indices that the model has not encountered during pre-training often generates numerous abnormal values, which result in out-of-distribution problems and hinder the convergence of fine-tuning~\cite{pi}. This challenge becomes even more severe when expanding the window from 4k to over $>$1000k, as this process entails creating more than 90\% novel positional indices. Second, effective fine-tuning generally depends on access to training sequences of sufficient length. At present, there are only a limited number of available datasets containing such extended sequences, especially those that surpass 1000k tokens. Furthermore, the computational cost associated with training on ultra-long texts is extremely high, demanding extensive computational time and significant GPU resources. Third, scaling the context window to extreme lengths leads to more diffuse attention distributions across a large number of token positions, ultimately resulting in reduced effectiveness on shorter context tasks~\cite{pi}.

A commonly employed strategy to address the initial challenge involves interpolating RoPE positional embeddings~\cite{rope,pi}, which rescales out-of-distribution positional indices to fit within the bounds established during pretraining, as depicted in Fig.\ref{fig:overview}. The Position Interpolation (PI) method~\cite{pi} achieves this by performing a linear interpolation of RoPE’s rotary angles according to the desired extension ratio. In contrast, NTK~\cite{ntk,dynamicntk} proposes an asymmetric interpolation and extrapolation across different RoPE dimensions. YaRN~\cite{yarn} introduces a grouping scheme, classifying RoPE dimensions into three distinct categories based on frequency, and then applying extrapolation, NTK-style adjustments, or linear interpolation to each group. Despite these efforts, the positional embedding in Transformer models is characterized by a \emph{complex, non-uniform distribution of information entropy}. Current methods do not fully exploit these nuanced irregularities, which results in a loss of positional information and consequently constrains the achievable context window.

Section~\ref{sec:analysis} presents two main empirical observations: \textbf{(1)} Robust positional interpolation necessitates addressing two sources of non-uniformity: the variability in RoPE dimensions and disparities among token positions. Reduced levels of interpolation are advantageous for lower RoPE dimensions as well as for tokens near the beginning of the sequence, although the ideal approach is influenced by the specific extended sequence length in question.
\textbf{(2)} When these non-uniform factors are incorporated into the interpolation process, the approach preserves more of the information contained in the original RoPE—especially with respect to crucial dimensions and token locations. This approach limits the information loss typically associated with positional interpolation, thereby enabling more effective initialization before fine-tuning. Additionally, it permits an up to 8$\times$ sequence length extension even when fine-tuning is not performed.

Building upon these insights, we present \textbf{LongRoPE}, a robust approach designed to increase the LLM context window to over 2 \emph{million} tokens. LongRoPE is underpinned by three primary innovations. Firstly, {LongRoPE} leverages multidimensional, non-uniform characteristics within positional interpolation to their fullest extent. It determines optimal rescale factors for the rotational angles in RoPE for every dimension, conditioned on specific token positions. Given that the search space for finding these rescale factors grows exponentially with the desired extension ratio, {LongRoPE} employs an evolutionary search algorithm, incorporating two specific optimization strategies to enhance the efficiency of this search process. An illustration of the resulting rescaled RoPE from this procedure is provided in Fig.~\ref{fig:overview}.

Subsequently, {LongRoPE} employs a resource-efficient, stepwise extension methodology that enables a context window of 2048k, circumventing the necessity to directly fine-tune models on ultra-long texts, which are both rare and difficult to procure. The approach starts by identifying an optimal 256k length for the pre-trained LLM, followed by fine-tuning at this sequence length. Owing to our non-uniform positional interpolation mechanism—which supports an 8$\times$ extension even without further fine-tuning—we next perform an additional search to determine suitable RoPE rescale factors using the fine-tuned, length-extended LLM. Through this process, the 2048k context window is attained for both LLaMA2 and Mistral~\cite{mistral}.

In order to prevent declines in performance when processing shorter input sequences, {LongRoPE} further tunes the RoPE rescale factors in the extended language model. Following the approach used for extending the context window from 256k to 2048k, the method also performs a downward rescaling to 4k and 8k context sizes on the LLM fine-tuned with a 256k window, utilizing the search algorithm to minimize positional interpolation effects. For inference, whenever the input sequence length falls below 8k tokens, the model applies the RoPE rescale factors determined during the search procedure.

Comprehensive evaluations conducted on multiple LLMs and a range of long-context tasks validate the efficacy of our approach. The results indicate that LongRoPE consistently maintains low perplexity across evaluation lengths spanning from 4k to 2048k, surpasses 90\% accuracy in passkey retrieval, and matches the performance of standard models on benchmarks set within a 4096 token context window. Furthermore, LongRoPE is compatible with all LLMs utilizing RoPE embeddings. We plan to publicly release both our code and the LongRoPE-2048k models.

\section{Non-uniformity in Positional Interpolation}
\label{sec:analysis}

\subsection{Preliminary}

Transformer architectures necessitate clear positional cues—typically achieved through position embeddings—to capture the sequential nature of input tokens. In this study, we investigate the RoPE~\cite{rope} position embedding, a method that has become prevalent in contemporary large language models (LLMs). For a token located at position index $n$, its RoPE encoding can be succinctly expressed as:

\begin{equation}
		\fontsize{7.8}{7.8} \selectfont
	\label{eq:rope}
	\left[ cos(n\theta_0), sin(n\theta_0), cos(n\theta_1), \cdot\cdot\cdot, cos(n\theta_{d/2-1}), sin(n\theta_{d/2-1}) \right]   
\end{equation}

where $d$ indicates the dimensionality of the embedding space, 
$n\theta_i$ denotes the rotary angle assigned to the token at the $n$-th position, 
and $\theta_i = \theta^{-2i/d}$ encodes the rotation frequencies. By default, RoPE sets the base value of $\theta$ to 10000.

\textbf{\emph{Context window extension ratio $s$} and positional interpolation}. The extension ratio $s$ is defined as the quotient of the new context window length $L'$ over the original context length $L$: $s = \frac{L'}{L}$.

In order to increase the context window from $L$ to $L'$, existing positional interpolation techniques propose modifying the rotation frequencies $\theta_i$ by scaling them down according to an extension ratio $s$. Defining $\beta=\theta^{2/d}$ and using $\lambda$ as the effective scaling parameter associated with $s$, we generalize these positional interpolation strategies as:

\begin{equation}	
	\fontsize{7.5}{7.5} \selectfont
	\label{eq:unified_rope}
	\begin{aligned}
		\left[cos\left(\frac{n}{\lambda(\beta)^0}\right), sin \left(\frac{n}{\lambda(\beta)^0}\right), cos\left(\frac{n}{\lambda(\beta)^1}\right), \cdot\cdot\cdot,  sin \left(\frac{n}{\lambda(\beta)^{d/2-1}}\right) \right]   
	\end{aligned}
\end{equation}

\textbf{Linear positional interpolation (PI)}. PI~\cite{pi} employs a linear interpolation approach for position indices when operating within the length range seen during pre-training. Given an intended sequence extension ratio $s$, the rotational angles for each position are uniformly scaled by a factor of $\lambda = s$ in all RoPE dimensions. This adjustment results in more compressed position representations, which impairs the model's capability to differentiate between tokens that are near each other in the sequence. Consequently, PI often yields suboptimal results as the extension ratio increases.

\textbf{NTK-based interpolation and extrapolation}. Works such as ~\cite{ntk,dynamicntk} analyze RoPE from the perspective of information encoding, utilizing Neural Tangent Kernel (NTK) theory~\cite{jacot2018neural,tancik2020fourier}. To address the problem of crowding at certain positions in PI, their approach aims to more evenly spread the interpolation demand across the various RoPE dimensions. This is achieved by applying less scaling to the lower (higher frequency) dimensions and greater scaling to the higher (lower frequency) dimensions, which enables both interpolation and extrapolation of positions where $\lambda=s^{i}$. The advanced dynamic NTK~\cite{dynamicntk} further refines this by dynamically setting the extension ratio for each position based on the actual sequence length. In contrast to PI, which relies on fine-tuning, NTK-aware approaches enable context window extension without the need for fine-tuning, although their maximum extension ratio typically does not exceed 4$\times$.

\textbf{YaRN}~\cite{yarn} presents a notable advancement in positional interpolation by segmenting RoPE dimensions into three distinct groups according to their frequencies, with each group assigned a specialized interpolation method. For dimensions with high frequencies, extrapolation is applied ($\lambda$=1); dimensions associated with low frequencies utilize linear interpolation (PI); and those in the intermediate frequency range employ the NTK approach. The core innovation of YaRN is this frequency-based grouping of RoPE dimensions. However, the assignment of these groups is currently reliant on manual, experience-driven experimentation, which could lead to less-than-optimal results when applied to newly developed LLMs.

\subsection{Study on Non-uniform Positional Interpolation}

Drawing motivation from NTK and YaRN, we observe that their improvements stem from introducing non-linearity—especially through treating various RoPE dimensions with distinct frequencies, thereby enabling tailored interpolation and extrapolation. Nevertheless, the prevailing forms of non-linearity are primarily guided by manually crafted heuristics. This leads us to inquire: (1) Does the existing approach to positional interpolation achieve optimality? (2) Might there be additional, yet-to-be-discovered non-linear transformations?

In order to address these issues, we employ evolution search (refer to Sec.~\ref{sec:method}) to identify improved non-uniform positional interpolation strategies for LLaMA2-7B. The evolution process is directed by perplexity scores, which are evaluated on 5 randomly selected samples from the PG19~\cite{pg19} validation set. Our experimental investigation yields several notable insights.

\textbf{Finding 1}: \textit{RoPE dimensions exhibit substantial non-uniformities, which are not effectively handled by current positional interpolation methods.}

We perform a search for the optimal $\lambda$ corresponding to each RoPE dimension as described in Eq.~\ref{eq:unified_rope}. Table~\ref{tbl:exp1} reports the perplexity outcomes of LLaMA2-7B, evaluated on the PG19 and Proof-pile~\cite{proof-pile} test datasets, across various methods and in the absence of fine-tuning. The solution identified through our search process yields substantial improvements, indicating that both existing linear (PI) and non-uniform (Dynamic-NTK and YaRN) interpolation strategies are not optimal. It is noteworthy that YaRN performs worse than both PI and NTK on PG19, failing to reach the intended context window for models without fine-tuning. Specifically, YaRN’s perplexity sharply increases beyond the 7k token mark within an 8k context window.

In our investigation, the rescaling coefficients $\lambda$ in Eq.~\ref{eq:unified_rope} exhibit variability, unlike the constant scaling factor $s$ utilized in PI, the formulation of NTK, and the group-based computation in YaRN. These variable scaling coefficients lead to notable gains in LLaMA2’s language modeling accuracy (as measured by perplexity) across 8k and 16k context lengths, all achieved without additional fine-tuning. The main reason for this enhancement is that the resultant positional embeddings largely retain the properties of the original RoPE—particularly in critical dimensions—thereby easing the model’s challenge in recognizing and differentiating among closely positioned tokens.

\textbf{Finding 2}: \textit{RoPE for the initial tokens in the input sequence should be extrapolated with less interpolation.} 

We propose that for the first $\hat{n}$ tokens in the input sequences, position interpolation by RoPE should be minimized, given that these tokens generally acquire higher attention weights and are thus more influential within attention layers, as documented in Streaming LLM~\cite{streamingllm} and LM-Infinite~\cite{han2023lminfinite}. To test this hypothesis, we extend the context window to 8k and 16k using both PI and NTK, while selectively excluding the initial $\hat n$ tokens (where $\hat n$ takes values from 0, 2, ..., to 256) from position interpolation. Setting $\hat n$=0 yields the standard PI and NTK setups. Table~\ref{tbl:tokenposition} presents two main findings: \textbf{(1)} Omitting position interpolation for the initial tokens results in measurable performance gains. \textbf{(2)} The best-performing value of $\hat n$ is contingent on the desired window extension length.

\textbf{Finding 3}: \textit{Non-uniform positional interpolation effectively extends LLM context window in both fine-tuning and non-fine-tuning settings.} 

Although our results demonstrate that the proposed non-uniform position interpolation method yields notable gains in extension capabilities at 8k and 16k contexts without the need for additional training, further increasing the context length necessitates fine-tuning. Accordingly, we perform fine-tuning on LLaMA2-7B utilizing our optimized RoPE configuration to accommodate a 64k context window (details available in the Appendix). As presented in Table~\ref{tbl:finetune}, our approach delivers substantial improvements over PI and YaRN, both prior to and following the fine-tuning of LLaMA2-7B. This advantage stems from the superior exploitation of non-uniform positional interpolation, which effectively reduces information loss and offers a more favorable starting point for the fine-tuning process.

\textbf{Summary}. This work identifies two critical sources of non-uniformity: differences across RoPE dimensions and disparities among token positions. By leveraging these forms of non-uniformity within positional interpolation, our approach achieves marked improvements in the context extension capabilities of LLMs.

\section{LongRoPE}

\label{sec:method}
Building upon these insights, we propose LongRoPE. This approach initially develops a computationally efficient search method designed to leverage both non-uniformities. Leveraging this, we subsequently extend the LLM context window to handle over 2 million tokens.

\subsection{Problem Formulation}

The presence of these two forms of non-uniformity significantly expands the range of potential solutions, thereby increasing the complexity of the optimization task. To manage this challenge, we conceptualize the multidimensional non-uniform position interpolation optimization as a search-based problem.

Consider a LLM configured for a context window size of $L'$ that is required to process long-form input documents $\mathbf{X}$, where each $\mathbf{x}\in\mathbf{X}$ exceeds $L'$ tokens in length. For the $i^{th}$ dimension in the RoPE embedding at token index $n$, the corresponding original rotary angle is given by $\frac{n}{\beta^i}$. The resulting optimization task can be described as follows:

\begin{equation}
\begin{gathered}
		\label{eq:problem}

\underset {\mathbf{x}\in\mathbf{X}; \, |\mathbf{x}|\ge L'}{ \operatorname {arg\,min} } \,  \mathcal{L}\, \left( \text{LLM} (\text{RoPE}, \mathbf{X})\right), \text{with} 
\\
\scalebox{0.93}{
$\underset {\begin{subarray}{l} i=0,\cdot\cdot,\frac{d}{2}-1;\\ n\in[ 0,|\mathbf{x}|);\end{subarray}}{\text{RoPE}_(n)}= \left[\cdots,\cos\left(\mathbb{I}(\hat{\lambda}_{i}, \hat{n}) \frac{n}{\beta^i}\right), \sin\left(\mathbb{I}(\hat{\lambda}_{i}, \hat{n}) \frac{n}{\beta^i}\right),\cdots\right]$}
\\
\text{where} \; \mathbb{I}(\hat{\lambda}_{i},\hat{n}) = \begin{cases}
1 & \text{if } n < \hat{n} \\
\frac{1}{\lambda_{i}} & \text{if } n \geq \hat{n}
\end{cases}
	\end{gathered}
\end{equation}

In this approach, we define a collection of scaling factors, $\mathbb{I}(\hat{\lambda}_{i},\hat{n})$, intended to address both types of non-uniformity present. Here, $\hat{\lambda}_i$ characterizes non-uniformities in the RoPE dimensions, while $\hat{n}$ specifies non-uniformity across token positions. The function $\mathbb{I}(\hat{\lambda}_{i},\hat{n})$ is employed to adjust the rotation angle associated with the $i^{th}$ RoPE dimension, where $\hat{\lambda}_i$ is used as the scaling factor, and $\hat{n}$ serves as the positional threshold for applying the adjustment. For the first $\hat{n}-1$ token positions, the scaling factor $\hat{\lambda}_i$ is not engaged, and the standard RoPE rotation angle $\frac{n}{\beta^i}$ remains in use. When $n\ge\hat{n}$, however, the scaling factor begins to influence the calculation.

Assuming a desired context window length of $L'$, our task is to determine the best set of rescale factors—$\mathbb{I}(\hat{\lambda}_{0},\hat{n})$, $\mathbb{I}(\hat{\lambda}_{1},\hat{n})$, ..., $\mathbb{I}(\hat{\lambda}_{i},\hat{n})$, and so forth—for each RoPE dimension from the first through the $d^{th}$. This adjustment enables the target $\text{LLM}$, utilizing the modified $\text{RoPE}$, to attain the lowest possible next-token prediction loss $\mathcal{L}$ (such as minimizing perplexity) on input data $\mathbf{X}$, where the token sequence is of length $L'$.

\subsection{Searching the Non-uniform Position Interpolation}

To address the issue outlined in Eq.~\ref{eq:problem}, we present a straightforward yet highly efficient approach that identifies the optimal RoPE rescale factors, thus capitalizing on the multidimensional, non-uniform nature of positional embeddings.

\textbf{Search space}. We construct an extensive search space to capture both types of non-uniformity. An overview of this space is provided in Table~\ref{tbl:searchspace}. More precisely, our method permits the assignment of a distinct rescale factor to each RoPE embedding dimension. For simplicity in the search space construction, we parameterize the search using $\lambda_i$ and $\hat{n}$, as opposed to directly optimizing $\mathbb{I}(\hat{\lambda}_{i},\hat{n})$, with the relationship $\hat{\lambda}_{i}=1/\lambda_i$. According to Table~\ref{tbl:searchspace}, each $\lambda_i$ may take values starting from a lower bound of 1.0 (corresponding to straightforward extrapolation) up to an upper bound of $s\times1.25$ (allowing for greater interpolation than PI), incremented in steps of 0.01. Here, $s$ denotes the desired context window extension ratio.

The parameter $\hat{n}$ specifies how many of the initial token positions maintain their original positional encoding without applying position interpolation (specifically, they utilize the standard RoPE embedding). In our experiments, we vary $\hat{n}$ over the set \{0, 1, 2, 4, 8, 12, 16, 20, 24, 28, 32, 64, 128, 256\}. Setting $\hat{n}=0$ means that rescale factors, as determined through searching, are applied to every token position.

\textbf{Evolution-based search}. The search space defined in Table~\ref{tbl:searchspace} encompasses a vast array of positional interpolation possibilities, making systematic exploration highly challenging. For instance, an extension factor of $s=4\times$ results in $400^{128/2}\times14=4\times10^{167}$ potential configurations. As the extension ratio increases, the number of possibilities grows exponentially. To efficiently navigate this enormous space, we employ evolution search~\cite{guo2020single}, supplemented by two optimization strategies that significantly enhance the efficiency of the search process. The complete search workflow is presented in Algorithm~\ref{alg:search}.

\textit{Optimized initial population generation}. In lieu of purely random initialization for the population of $P$ rescale factors, we explicitly include the RoPE rescale factors associated with PI, NTK, and YaRN as distinct individuals at the outset. The rest of the initial population, amounting to $P-3$ individuals, is formed by randomly mutating these three reference rescale factors with probability $p$.

\textit{Monotonically non-decreasing constraint}. Following the creation of the initial population, we evaluate the LLM perplexity for every individual. This involves applying the associated RoPE rescale factors to the target LLM and then determining the perplexity of input $\mathbf{X}$. The best-performing $k$ individuals are selected as parents for subsequent generations. Nevertheless, the expansive search space may result in unproductive mutations and crossovers, which can direct the process toward suboptimal regions and trigger excessive perplexity evaluations. This inefficiency becomes more pronounced as $L'$ increases, since each perplexity calculation requires substantial inference time.

To mitigate this issue, we enforce a non-decreasing monotonicity restriction on the chosen RoPE rescaling coefficients, such that $\lambda_i\le \lambda_{i+1}$. Only those RoPE configurations adhering to this monotonicity condition are utilized in the LLM perplexity assessments, which considerably decreases the computational burden of the search. More precisely, we stipulate that the sequence $\lambda_i$ must monotonically increase with respect to the RoPE dimension (for $i = 0, \ldots, 63$). This requirement for monotonicity along the dimension is informed by NTK theory~\cite{jacot2018neural,tancik2020fourier,ntk}, which posits that dimensions with higher frequency (corresponding to lower indices) necessitate less interpolation—thus a smaller $\lambda_i$—while dimensions with lower frequency (higher indices) can accommodate greater interpolation, necessitating a larger $\lambda_i$.

\begin{algorithm}[t]
	\small
	\caption{The search algorithm}
	\textbf{Input:} target LLM, input samples $\mathbf{X}$,   population size $P$, mutation size $N_1$, crossover size $N_2$, max iterations $\mathcal{T}$, mutate probability $p$\\

	\begin{algorithmic}[1]
		\label{alg:search}
		\STATE Initialize Top-k as empty set;	
		\STATE $\text{P}_0$ = \textit{Initialize\_population\_with\_optimization}($P$, $\mathbf{X}$, $p$); 
		\FOR{$i$ from $1$ to $\mathcal{T}$}
		\STATE \textit{Compute\_perplexity} (LLM, $\text{P}_{i-1}$, $\mathbf{X}$);
		\STATE Top-k = \textit{Update\_Topk}(Top-k, $\text{P}_{i-1}$);
		\STATE $\text{P}_{mutation}$ = \textit{Mutation\_with\_mono\_constraint}(Top-k, $N_1$, $p$); 
		\STATE $\text{P}_{crossover}$ = \textit{Crossover\_with\_mono\_constraint}(Top-k, $N_2$);
		\STATE $\text{P}_i$ = $\text{P}_{mutation}$ $\cup$ $\text{P}_{crossover}$ $\cup$ Top-k;
		\ENDFOR
		\STATE Output the member of Top-k achieving the minimum perplexity;
	\end{algorithmic}
\end{algorithm}

\textbf{8$\times$ extension without fine-tuning}. Through evolutionary search, our approach efficiently selects non-uniform RoPE rescale factors that safeguard crucial dimensions and positions, thereby reducing information loss from interpolation. As illustrated in Fig.\ref{fig:non-ft}, this strategy enables LLaMA2’s context window to expand from 4k up to 32k tokens, all without the need for fine-tuning. By comparison, alternative approaches like PI, as well as non-uniform NTK and YaRN, result in a significant increase in perplexity once the window exceeds a 2$\times$ extension.

\subsection{Extending LLM Context Window to 2048K}
\label{sec:extendto2048k}

\textbf{Progressive extension to 2048k}. In this section, we present our approach for scaling the context window of pre-trained LLMs from the conventional 4k up to more than 2048k tokens. Our results indicate that utilizing non-uniform positional interpolation permits an 8$\times$ increase without the need for fine-tuning. However, for more substantial extensions, such as a 512$\times$ increase, fine-tuning becomes essential. One strategy involves determining suitable RoPE rescale factors at the desired 2048k context length prior to initiating fine-tuning. Nevertheless, this approach is hampered by the immense computational costs required for training. Furthermore, our observations suggest that effectively fine-tuning LLMs at such large extension scales poses considerable difficulties (refer to Appendix).

Fortunately, LongRoPE demonstrates efficacy in both the base model and the extended LLM after fine-tuning. Accordingly, we propose an efficient and incremental approach that reaches the desired 2048k target in only 1,000 fine-tuning steps, requiring a maximum training sequence length of 256k.

$\Diamond$ \textit{LongRoPE search for expanding pre-trained LLM to 256k}. Using LLaMA2 as a case study, we perform a search procedure targeting context window lengths of 128k and 256k. In this phase, the window is enlarged by factors of 32$\times$ and 64$\times$, respectively.

$\Diamond$ \textit{Fine-tuning to 256k}. Next, the pre-trained LLM undergoes fine-tuning to extend the context window to 256k. In detail, LLaMA2 is initially fine-tuned for 400 steps using the RoPE rescaling factors calibrated for 128k. Afterwards, these factors are switched to those appropriate for 256k on the obtained checkpoint, followed by another 600 steps of fine-tuning. This staged approach demonstrates greater efficiency compared to immediate fine-tuning at 256k.

$\Diamond$ \textit{Scaling fine-tuned extended LLM to 2048k via LongRoPE search}. In the concluding stage, we apply a subsequent search to the LLM that has already been fine-tuned for a 256k context length. This approach enables the attainment of a significantly expanded context window of 2048k tokens, all without additional rounds of fine-tuning. The resulting overall extension factor amounts to $512\times$.

\textbf{Performance decline within original context window}. When expanding the context window to a substantial length of 2048k, we observe a deterioration in performance for inputs confined to the initial context length. This phenomenon has been documented with positional interpolation~\cite{pi}, where the embedding of positions in higher dimensions within the original context becomes compressed into a more constrained space, thereby impairing the language model's efficacy. Specifically, applying a 512$\times$ extension ratio results in substantial positional congestion in the original 4k context window.

To address this, an additional evolution search is conducted on the expanded LLM to fine-tune the RoPE rescale factors specifically for shorter context lengths, such as 4k and 8k. The upper bound for the searched $\lambda$ is decreased, since shorter lengths necessitate less positional interpolation. At inference time, the LLM automatically modifies the relevant RoPE rescale factors.

\section{LongRoPE}

\label{sec:method}
Building on these observations, we propose LongRoPE. This method initially applies an optimized search algorithm designed to leverage both types of non-uniformity, subsequently employing this strategy to push the LLM context window past the 2 million token threshold.

\subsection{Problem Formulation}

The presence of these two non-uniformities significantly expands the solution space and complicates the optimization process. To manage this, we reformulate the multidimensional non-uniform position interpolation optimization challenge as a search problem.

Consider an LLM designed for a context window of size $L'$ and tasked with processing long input documents $\mathbf{X}$, where each element $\mathbf{x}\in\mathbf{X}$ contains more tokens than $L'$. Let the initial rotary angle assigned to the $i^{th}$ dimension within the RoPE embedding at token index $n$ be represented by $\frac{n}{\beta^i}$. With this setup, the corresponding optimization objective is defined as:

\begin{equation}
\begin{gathered}
		\label{eq:problem}

\underset {\mathbf{x}\in\mathbf{X}; \, |\mathbf{x}|\ge L'}{ \operatorname {arg\,min} } \,  \mathcal{L}\, \left( \text{LLM} (\text{RoPE}, \mathbf{X})\right),	\text{with \;} 
\\
\scalebox{0.93}{
$\underset {\begin{subarray}{l} i=0,\cdot\cdot,\frac{d}{2}-1;\\ n\in[ 0,|\mathbf{x}|);\end{subarray}}{\text{RoPE}_(n)}= {\left[\cdots, \cos\left(\mathbb{I}(\hat{\lambda}_{i}, \hat{n})\frac{n}{\beta^i}\right),  \sin\left(\mathbb{I}(\hat{\lambda}_{i}, \hat{n})\frac{n}{\beta^i}\right),\cdots\right] }$}\\
	\text{where \;} \mathbb{I}(\hat{\lambda}_{i},\hat{n})=\begin{cases}
	1&\mbox{\;  $n<\hat{n}$}\\
{\frac{1}{\lambda}_{i}}&\mbox{\; $n\geq\hat{n}$}
\end{cases}
	\end{gathered}
\end{equation}

In this context, we define a group of rescaling factors, $\mathbb{I}(\hat{\lambda}_{i},\hat{n})$, to address both types of non-uniformities present. Here, $\hat{\lambda_i}$ corresponds to irregularities in RoPE dimensions, while $\hat{n}$ refers to the irregularity associated with token positions. We employ $\mathbb{I}(\hat{\lambda}_{i},\hat{n})$ as a means to adjust the rotation angle along the $i^{th}$ RoPE dimension; $\hat{\lambda}_{i}$ serves as the dimension-specific rescale factor and $\hat{n}$ designates a positional threshold. For the first $\hat{n}-1$ tokens, the rescaling is inactive and the standard RoPE rotation angle $\frac{n}{\beta^i}$ is maintained. Once the token index reaches $n\ge\hat{n}$, the rescale factor begins to modify the rotation angle.

Suppose we have a desired context window length of $L'$. Our aim is to determine the set of optimal rescaling coefficients ($\mathbb{I}(\hat{\lambda}_{0},\hat{n})$, $\mathbb{I}(\hat{\lambda}_{1},\hat{n})$, ..., $\mathbb{I}(\hat{\lambda}_{i},\hat{n})$, ...) across all $d$ RoPE dimensions, starting from the $1^{st}$ to the $d^{th}$. By applying these rescaled values to the target $\text{LLM}$'s $\text{RoPE}$ module, we seek to minimize the next-token prediction loss, $\mathcal{L}$ (i.e., perplexity), evaluated on input samples $\mathbf{X}$ whose sequence length is $L'$.

\subsection{Searching the Non-uniform Position Interpolation}

To address the challenge presented in Eq.~\ref{eq:problem}, we propose a straightforward yet powerful approach that identifies optimal RoPE rescale factors, leveraging the multidimensional heterogeneity inherent in position embeddings to the fullest extent.

\textbf{Search space}. We construct an expansive search space that comprehensively accounts for both types of non-uniformity. The configuration of this search space is detailed in Table~\ref{tbl:searchspace}. In particular, we permit each dimension within the RoPE embedding to be assigned its own distinct rescale factor. To streamline the process of defining the search space, our approach involves searching over the parameters $\lambda_i$ and $\hat{n}$, instead of directly searching for $\mathbb{I}(\hat{\lambda}_{i},\hat{n})$, with $\hat{\lambda}_{i}=1/\lambda_i$. According to Table~\ref{tbl:searchspace}, the values of $\lambda_i$ are selected from a range starting at 1.0 (representing direct extrapolation) up to $s\times1.25$ (indicating a greater degree of interpolation than what Position Interpolation achieves), incremented by 0.01. Here, $s$ corresponds to the targeted expansion ratio for the context window.

The parameter $\hat{n}$ determines how many of the initial token positions preserve their original representations, specifically by maintaining the original RoPE embedding rather than applying position interpolation. In practice, $\hat{n}$ is selected from the set \{0, 1, 2, 4, 8, 12, 16, 20, 24, 28, 32, 64, 128, 256\}. Setting $\hat{n}=0$ indicates that all token positions are subject to the rescale factors identified during the search process.

\textbf{Evolution-based search}. The search space outlined in Table~\ref{tbl:searchspace} encompasses a vast array of positional interpolation configurations, making exhaustive exploration highly impractical. For instance, with an extension factor of $s=4\times$, there are $400^{128/2}\times14$=4$\times10^{167}$ possible combinations. As the extension ratio increases, the size of the search space grows at an exponential rate. To tackle this challenge, we employ evolution search~\cite{guo2020single} alongside two specialized optimization strategies that substantially enhance the efficiency of the search. The overall process is depicted in Algorithm~\ref{alg:search}.

\textit{Optimized initial population generation}. Rather than relying solely on random initialization for the population of $P$ rescale factors, we directly incorporate the three RoPE rescale factors associated with PI, NTK, and YaRN as members of the initial set. The other $P$-3 individuals are produced by applying random mutations to these three base rescale factors, each mutation occurring with probability $p$.

\textit{Monotonically non-decreasing constraint}. Following the initialization of the population, we evaluate the LLM perplexity for every individual. This involves applying the assigned RoPE rescale factors within the target LLM and measuring the perplexity on the input $\mathbf{X}$. The highest-performing $k$ individuals are then selected as parents for subsequent evolutionary steps. Nevertheless, the enormity of the search space means that unrefined mutation and crossover processes can frequently result in suboptimal candidates, which incurs redundant perplexity assessments. This inefficiency becomes more pronounced as $L'$ grows, given the high computational cost associated with each perplexity evaluation.

To tackle this issue, we enforce a non-decreasing monotonicity requirement on the selected RoPE rescaling parameters: $\lambda_i \leq \lambda_{i+1}$. Only those RoPE configurations that meet this criterion are utilized in the LLM for perplexity measurement, thereby greatly curbing the computational overhead of the search. More precisely, we mandate that $\lambda_i$ is monotonically increasing as the RoPE dimension rises (for $i = 0, \ldots, 63$). This constraint on dimension-wise monotonicity aligns with principles from NTK theory~\cite{jacot2018neural,tancik2020fourier,ntk}, which indicates that the lower dimensions—associated with higher frequencies—necessitate less interpolation (smaller $\lambda_i$), while the higher dimensions—associated with lower frequencies—can tolerate or benefit from more interpolation (larger $\lambda_i$).

\begin{algorithm}[t]
	\small
	\caption{The search algorithm}
	\textbf{Input:} target LLM, input samples $\mathbf{X}$, population size $P$, mutation size $N_1$, crossover size $N_2$, maximum number of iterations $\mathcal{T}$, mutation probability $p$\\

	\begin{algorithmic}[1]
		\label{alg:search}
		\STATE Initialize Top-k as empty set;	
		\STATE $\text{P}_0$ = \textit{Initialize\_population\_with\_optimization}($P$, $\mathbf{X}$, $p$); 
		\FOR{$i$ from $1$ to $\mathcal{T}$}
		\STATE \textit{Compute\_perplexity} (LLM, $\text{P}_{i-1}$, $\mathbf{X}$);
		\STATE  Top-k $\gets$ \textit{Update\_Topk} (Top-k, $\text{P}_{i-1}$);
		\STATE $\text{P}_{mutation}$ $\gets$ \textit{Mutation\_with\_mono\_constraint} (Top-k, $N_1$, $p$); 
		\STATE $\text{P}_{crossover}$ $\gets$ \textit{Crossover\_with\_mono\_constraint} (Top-k, $N_2$);
		\STATE $\text{P}_i$ = $\text{P}_{mutation}$ $\cup$ $\text{P}_{crossover}$ $\cup$ Top-k;
		\ENDFOR
		\STATE Output the member of Top-k that achieves the minimal perplexity;
	\end{algorithmic}
\end{algorithm}

\textbf{8$\times$ extension without fine-tuning}. By utilizing evolutionary search, we successfully determine optimal non-uniform RoPE rescale factors that maintain essential dimensions and positional encodings, thereby reducing information loss caused by interpolation. As illustrated in Fig.\ref{fig:non-ft}, this approach enables LLaMA2’s context window to be expanded from 4k to 32k without necessitating any fine-tuning. Conversely, alternative techniques like PI and non-uniform variants of NTK and YaRN exhibit significant increases in perplexity following only a 2$\times$ context window extension.

\subsection{Extending LLM Context Window to 2048K}
\label{sec:extendto2048k}

\textbf{Progressive extension to 2048k}. In this section, we present our approach for scaling the context window of pre-trained LLMs from the conventional 4k up to beyond 2048k. Our results indicate that utilizing non-uniform positional interpolation enables an 8$\times$ increase in context length without the need for fine-tuning. However, achieving substantially larger extensions (such as 512$\times$) necessitates fine-tuning. One possible strategy is to identify appropriate RoPE rescaling factors corresponding to the desired 2048k context size, followed by fine-tuning. Nevertheless, this approach is hampered by the considerable computational resources required. Additionally, our empirical observations reveal that effective fine-tuning becomes increasingly difficult at such high extension ratios (refer to Appendix for further details).

Fortunately, LongRoPE demonstrates efficacy with both base and extended fine-tuned LLMs. Accordingly, we propose a streamlined, incremental approach that attains the desired 2048k target using only 1k fine-tuning iterations and confines training to a maximum sequence length of 256k.

$\Diamond$ \textit{Scaling pre-trained LLMs to 256k via LongRoPE search}. Using LLaMA2 as a representative model, we perform a search aiming for context window lengths of 128k and 256k. The context window is thus expanded by factors of 32$\times$ and 64$\times$ at this stage.

$\Diamond$ \textit{Fine-tuning to 256k}. In this phase, we adapt the pre-trained LLM to handle a context window of 256k tokens. Our approach involves initially fine-tuning LLaMA2 for 400 steps utilizing RoPE rescaling parameters corresponding to 128k. Following this, we substitute these parameters with those suited for 256k and proceed with another 600 steps of fine-tuning starting from the updated checkpoint. This staged procedure demonstrates greater efficiency compared to commencing the fine-tuning directly at 256k.

$\Diamond$ \textit{LongRoPE search for 2048k extension of fine-tuned LLM}. In the final stage, we apply a second round of search on the LLM already fine-tuned to a 256k sequence length. Through this approach, we are able to achieve an exceptionally large context window of 2048k, all without requiring additional fine-tuning steps. This yields an overall extension factor of $512\times$.

\textbf{Recovery of performance in shorter context windows}.  When expanding the context window to a vast length of 2048k, we observe a degradation in performance inside the initial context range. This drawback has been previously reported for positional interpolation~\cite{pi}, as it compels the positional embeddings that fall within the original context window to be compressed into a more confined subspace, which diminishes the efficacy of the language model. In the case of a 512$\times$ expansion, the original 4k context window's positions are subject to even greater concentration.

To address this issue, we conduct an additional evolution-based search on the augmented LLM to fine-tune RoPE rescale factors specifically for shorter context windows (such as 4k and 8k). The upper bound for the searched $\lambda$ is lowered because shorter sequences necessitate less positional interpolation. At inference time, the LLM adaptively modifies the relevant RoPE rescale factors as needed.

 \section{Experiments}

\label{sec:exp}
\subsection{Setup}
\textbf{Evaluation Tasks and Models}. We incorporate LongRoPE into LLaMA2-7B and Mistral-7B, assessing their capabilities based on three criteria: (1) perplexity scores for language models with extended context lengths when processing lengthy documents; (2) the Passkey retrieval task, designed to test how effectively a model can identify a straightforward passkey amidst a large volume of distractor text; and (3) performance on conventional LLM benchmarks, using a standard context window of 4096 tokens.

\textbf{Fine-tuning}. In the case of LLaMA2, a linear learning rate decay schedule is employed, starting at 2e-5, and a global batch size of 32 is set. Fine-tuning is performed over 400 iterations using the Redpajama~\cite{redpajama} dataset, which is segmented into 128k token-long sequences, each demarcated by BOS and EOS tokens. Subsequently, building on the resulting checkpoint, the model undergoes a further 600 steps of training to extend the context window to 256k. The initial 128k context configuration utilizes 8 A100 GPUs managed by the distributed training framework~\cite{cube}, whereas expanding to 256k context necessitates 16 A100 GPUs. For Mistral, a fixed learning rate of 1e-6 is applied, with a global batch size of 64. For both 128k and 256k configurations, training protocols are aligned with those outlined in YaRN~\cite{yarn}, running for 400 steps on Together Computer’s Long-Data Collections~\cite{mistraldata} using sequences of 16k tokens, and utilizing 4 A100 GPUs during training.

\textbf{Search}. When the target window size does not exceed 256k, the following configuration is applied: $P$=64, $N_1$=$N_2$=16, $p$=0.3, and $\mathcal{T}$=40. During each search iteration, the top 32 candidates are chosen for mutation and crossover procedures. Perplexity evaluation utilizes 5 randomly selected validation samples from PG19, ensuring that each sample meets the minimum target context length. For window sizes above 512k, all population, mutation, and crossover parameters are reduced by half. In these larger settings, perplexity is instead assessed on 3 randomly chosen validation samples drawn from Pile-Books3~\cite{pile}.

\textbf{Baselines}. To extend the context window to 2048k, we performed fine-tuning on models using 128k and 256k token contexts. As a result, we obtain LongRoPE-2048k (ft=128k) for LLaMA2 and LongRoPE-2048k (ft=256k) for Mistral. These four models are evaluated against leading methods for increasing context length, particularly focusing on open-source LLMs that have been fine-tuned with positional interpolation techniques such as PI, NTK, and YaRN. The comparison set comprises Together-32k~\cite{together}, Code LLaMA~\cite{codellama}, LongLoRA-full-FT-100k~\cite{longlora}, YaRN-LLaMA, and YaRN-Mistral~\cite{yarn}.

\subsection{Main Results}

\label{sec:mainresult}
\textbf{Long sequence language modeling at 256k tokens}. Our initial evaluation focuses on benchmarking against leading long-context LLMs on sequences up to 256k tokens. To highlight the robustness of our approach, we conduct experiments on two datasets: Proof-pile~\cite{pg19} and the PG19~\cite{pile} test sets. We assess perplexity across a range of context sizes using a 256-token sliding window. For the PG19 dataset, all 100 test documents are included in the analysis. For Proof-pile, in alignment with YaRN~\cite{yarn}, we randomly select 10 examples, each comprising at least 128k tokens.

Table~\ref{tbl:proof-pile} and Table~\ref{tbl:pg19} present perplexity results for LLaMA2 and Mistral, each adapted using various interpolation techniques on the Proof-pile and PG19 benchmarks. From these results, two main points emerge: \textbf{(1)} our models extended to longer context windows consistently achieve lower perplexity as evaluation length increases from 4k up to 256k, indicating effective exploitation of extended context. \textbf{(2)} Despite operating with context windows up to 16 times longer than standard settings—a scenario that typically hinders short-length performance—our LongRoPE-2048k models surpass leading baseline models within the 256k context window.

\textbf{Long sequence language modeling beyond 2000k}. To assess performance on very lengthy texts, we utilize the Books3~\cite{pile} dataset. For efficient evaluation, a random sample of 20 books—each with a length greater than 2048k—is chosen, employing a sliding window approach with a size of 256k.

As indicated in Table~\ref{tbl:books3}, LongRoPE is able to extend the context window of both LLaMA2-7B and Mistral-7B up to 2048k, while maintaining perplexity that matches or surpasses baseline results at shorter sequence lengths ranging from 8k to 128k. Notably, there are substantial disparities in performance between LLaMA2 and Mistral at the 2048k context length. While Mistral exceeds the performance of baselines on shorter sequences, its perplexity rises above 7 for sequence lengths beyond 256k. In the case of LLaMA2, the model’s perplexity generally decreases as context length increases, though slight increases are observed at 1024k and 2048k. Additionally, for LLaMA2, using LongRoPE-2048k with a 256k fine-tuning length yields superior results compared to 128k, attributable to the smaller secondary extension ratio (i.e., 8$\times$ as opposed to 16$\times$). Conversely, Mistral exhibits better outcomes when fine-tuned with a window size of 128k. This difference is primarily due to employing a 16k training length for both 128k and 256k fine-tuning on Mistral, in line with YaRN’s configuration, which impacts Mistral's capability to further scale its context window following fine-tuning.

\textbf{Passkey retrieval}. Next, we investigate the influence of context window size on performance in generation settings. To do this, we adopt the passkey retrieval synthetic evaluation introduced by ~\cite{passkey}. In this scenario, the model's objective is to identify a random passkey (specifically, a five-digit numeral) embedded somewhere within an extended document. Details regarding the prompt template can be found in the appendix. For our assessment, we conduct 10 runs of the retrieval task, with the passkey positioned at randomly selected points uniformly distributed throughout the evaluation context window.

Fig.~\ref{fig:passkey} presents a comparative analysis of retrieval accuracy against several baselines. The performance of prior models deteriorates sharply, with their accuracy falling to zero once the context exceeds 128k tokens. By contrast, LongRoPE-LLaMA2-2048k (ft=256k) demonstrates strong robustness, maintaining a retrieval accuracy of at least 90\% across a wide range from 4k to 2048k tokens, even under the extremely demanding conditions of extracting a passkey from over a million tokens. Meanwhile, LongRoPE-Mistral-2048k (ft=128k) achieves perfect accuracy up to 1800k tokens, only decreasing to 60\% at the 2048k mark. This result corresponds with the trends observed in Table~\ref{tbl:books3}, where a moderate increase in perplexity is evident at 2048k tokens.

\textbf{Evaluation on standard benchmarks with the native context length}. The LongRoPE-2048k models are assessed on their native context window employing the Hugging Face Open LLM Leaderboard~\cite{huggingfaceboard} across both zero-shot and few-shot scenarios. Specifically, we conduct experiments with 25-shot ARC-Challenge~\cite{allenai:arc}, 10-shot HellaSwag~\cite{zellers2019hellaswag}, 5-shot MMLU~\cite{mmlu}, and 0-shot TruthfulQA~\cite{lin2021truthfulqa}.

As presented in Table~\ref{tbl:commonsense}, our models perform on par with the original benchmark, which was intended for shorter context windows, and surpass the original Mistral on TruthfulQA by a margin of +0.5\%. LongRoPE-LLaMA2-2048k, which was fine-tuned using a 256k context window, exhibits a minor drop in performance; nevertheless, its results remain largely acceptable across the majority of tasks.

\subsection{Ablation Results}

\textbf{Effectiveness of the second positional interpolation}. As part of our progressive extension method, we employ our search algorithm to perform a secondary non-uniform positional interpolation on the already fine-tuned, extended LLMs. To assess its efficacy, we carry out experiments with the fine-tuned LLaMA2-256k model. We further expand this model to sequence lengths of 512k, 1024k, and 2048k using both PI and YaRN for comparison. According to Table~\ref{tbl:secondsearch}, our approach with non-uniform positional interpolation consistently maintains perplexity at a stable level. By contrast, both PI and YaRN exhibit rapidly increasing perplexity as the extension ratio grows.

\textbf{Effectiveness of recovery at shorter context lengths.} To address the decline in performance observed at reduced context lengths, we fine-tune the RoPE factors for LongRoPE-2048k using our search strategy. In particular, the maximum permissible scale factors during the search are lowered to favor reduced interpolation when handling shorter 4k and 8k contexts. Table~\ref{tbl:shortperformance} presents a perplexity evaluation for LongRoPE-LLaMA2-2048k on Proof-pile at the 4k and 8k levels, alongside average accuracy across various LLM benchmarks. These findings highlight substantial gains in model performance at shorter context ranges.

\textbf{Analysis on the two forms of non-uniformities}. To assess the individual impact of each non-uniformity component, we conduct an ablation study. Specifically, we design two experimental settings: (i) scaling LLaMA2-7B to sequence lengths of 16k and 32k with varying strategies—using PI, optimizing only the RoPE dimension, and optimizing both non-uniformities; (ii) extrapolating a 256k fine-tuned LLaMA2 model up to 2048k in the same manner. Perplexity is measured without further fine-tuning. As shown in Table~\ref{tbl:searchdimension}, altering the RoPE dimension's non-uniformity yields a substantial decrease in perplexity compared to the linear interpolation approach of PI. Adjusting non-uniformity in the token position demonstrates clear performance gains at lengths of 16k and 32k, although this benefit does not persist at 2048k, potentially due to the extreme sequence length. Simply retaining original tokens without interpolation proves ineffective in these conditions, which we suggest as a direction for future investigation.

\section{Related Works}

Beyond position interpolation-based methods, this section examines alternative strategies from related literature.

\textbf{Retrieval-based} approaches leverage an external memory module to store extensive historical context and employ retrieval mechanisms to access pertinent documents during inference~\cite{longllama,wang2023augmenting,borgeaud2022improving}. Such frameworks often require substantial changes to the underlying LLM architectures. In comparison, our method is comparatively lightweight, requiring only minimal adjustments to positional embeddings. Additionally, our approach accommodates a broader range of long context scenarios, including tasks like long document summarization and few-shot learning, which extend beyond the scope of retrieval.

\textbf{Attention-based context window extensions}. In addition to interpolating positional embeddings, several studies have explored expanding the effective context window of LLMs by directly adjusting attention mechanisms, while keeping the original context window length unchanged~\cite{han2023lminfinite, streamingllm,parallelwindow}. Central to these approaches is the use of innovative attention masking techniques, which address the problem of escalating attention complexity introduced by additional positions. These approaches can be used in tandem with positional interpolation strategies, as they address orthogonal aspects of the problem.

\textbf{Fine-tuning based approaches} center on strategies to fine-tune pre-trained LLMs by adapting their position embeddings to handle longer contexts. Prior research such as Code LLaMA~\cite{codellama}, LLaMA2 Long~\cite{xiong2023effective}, and ScaledRoPE~\cite{liu2023scaling} adopts an approach where a significantly increased base value for RoPE is selected, followed by fine-tuning at the intended sequence length. In contrast, our approach provides adaptability across diverse target lengths and is capable of supporting sequences exceeding 2M tokens. Recently, due to the heavy GPU demands of fine-tuning models for extensive context windows (greater than 128k tokens), methods such as LongLoRA~\cite{longlora} and PoSE~\cite{zhu2023pose} have been introduced to alleviate these computational costs. Our technique complements these efficient fine-tuning methods and can be combined with them.

\section{Conclusion}

In this study, we introduce LongRoPE, a technique that significantly increases the context window of LLMs to an unparalleled 2048k, while preserving their performance on tasks requiring the standard, shorter contexts. By leveraging two distinct non-uniform properties of the RoPE positional embedding and employing an efficient evolutionary search strategy, our approach achieves two major advantages: it serves as an effective initialization for subsequent fine-tuning and allows for an 8$\times$ context expansion even in the absence of fine-tuning. Furthermore, we develop a stepwise extension methodology, utilizing LLMs fine-tuned on 256k-length sequences to progressively scale to a 2048k context window, again without additional fine-tuning. Comprehensive experimental results substantiate the efficacy of LongRoPE. We anticipate that LongRoPE-2048k will facilitate the development of numerous novel long-context applications and motivate future investigation in this area.

\section*{Broader Impacts} The research detailed in this paper seeks to contribute to progress in the domain of Machine Learning. While our work may have various implications for society, we do not identify any particular issues that warrant special emphasis at this time.

	\balance

\end{document}